\chapter{Regimen Rendiendi}

\epigraf{L'obsessió per la mesura destrueix l'autenticitat del procés.}{Byung-Chul Han}

\lettrine[loversize=0.2,lines=2]{L}{a} cultura de l'avaluació ha instaurat una paradoxa essencial: l'a-\ prenentatge, en lloc de ser un acte lliure d’exploració i descoberta, ha quedat sotmès al judici permanent de les notes, els resultats i les mètriques. Aquesta dinàmica converteix l’esforç intel·lectual en una persecució compulsiva d'una perfecció il·lusòria, desviant l’aten-\ ció de l'aprenentatge autèntic cap a la satisfacció d’una expectativa externa.

Quan l’aprenentatge depèn principalment de l’avaluació, deixa de ser un fi en si mateix per convertir-se en un mitjà. Els alumnes no cerquen comprendre, sinó demostrar. Això implica una inversió profunda en el sentit original del coneixement, que era entendre per créixer, i no entendre per complir. L’educació, sota aquesta lògica, es converteix en un cercle viciós on l’objectiu últim és l’aprovació externa, i no el desenvolupament intern.

Aquest perfeccionisme induït genera inevitablement angoixa. La ment que aprèn sota la pressió constant de ser avaluada no experimenta la curiositat com un goig natural, sinó com un risc. Aprendre deixa de ser descobrir per convertir-se en justificar-se. El pensament autèntic queda relegat a un segon pla, substituït per la repetició superficial que només busca garantir la valoració positiva.

La conseqüència d’aquesta pressió incessant és una espiral que consumeix l’energia mental, però que paradoxalment no aporta profunditat real. L’excés de perfeccionisme no només no millora la qualitat del que es fa, sinó que la redueix, atrapant l’individu en un cercle obsessiu on cada esforç es torna una reiteració buida del mateix acte de complaença.

Però l’autèntic aprenentatge no pot ser un acte d’ansietat, sinó un acte de llibertat. Aprendre implica assumir la imperfecció, la provisionalitat i l’error com a condicions necessàries del procés. Sense aquest reconeixement no hi pot haver un coneixement profund ni genuí, perquè només des de la llibertat de poder equivocar-se s’arriba a entendre realment.

Per superar aquesta dinàmica cal tornar a posar l’aprenentatge com a centre absolut, per damunt de qualsevol judici extern. L’avalua-\ ció hauria de ser tan sols un reflex, un instrument puntual, no un principi rector. Només així es pot recuperar el sentit original de la formació: no acumular puntuacions, sinó desenvolupar capacitats.

La llibertat d’aprendre sense la càrrega permanent del judici numèric permet redescobrir el plaer intrínsec de comprendre. És llavors quan es produeix l’aprenentatge autèntic, quan el coneixement deixa de ser una mera exigència per convertir-se en una forma profunda i vital de ser al món.

Tanmateix, aquesta reconfiguració exigeix més que una simple renúncia a l’obsessió per la nota: requereix transformar l’espai educatiu en un àmbit de confiança mútua on l’error no sigui càstig, sinó llavor de curiositat. Això reclama un canvi de paradigma en el qual el professor deixi de ser jutge per esdevenir guia i co-investigador. Cal, doncs, establir mecanismes d’avaluació regenerativa, basats en el diàleg obert i la retroalimentació qualitativa, que posin l’accent en el procés i no pas en el producte final.

A més, cal redefinir l’èxit: allunyar-nos de l’excel·lència quantitativa i apropar-nos a l’excel·lència relacional. L’objectiu serà no tant acumular mèrits, sinó construir comunitats d’aprenentatge on la curiositat sigui compartida i la responsabilitat, col·lectiva. L’èxit es mesurarà per la capacitat d’obrir preguntes, de sostenir dubtes i de teixir vincles intel·lectuals, més que per l’índex d’aprovat o suspès.

Finalment, recuperar l’autenticitat de l’aprenentatge implica constituir espais de reflexió lliure, fora de la dinàmica frenètica de l’examen, on l’estudiant trobi temps per aturar-se, qüestionar-se i imaginar noves potències. Només així, dotant d’espais i temps el jaç de la curiositat, es podrà revertir l’ombra que projecta el sistema de mesura sobredimensionada i tornar a l’horitzó original: l’aprenentatge com a acte de creació i de llibertat.